\documentclass{situationist-preprint}

\title{The Threshold of Restriction}
\msid{POL-16-b}
\msfamily{POL-16}
\msdiscipline{Political Philosophy}
\mstype{Companion manuscript}
\mskeywords{toleration; intolerance; militant democracy; threat assessment; Popper; Rawls}

\begin{document}
\maketitle

\begin{abstract}
POL-16-a argues that a society should tolerate all views except active intolerance, and should restrict active intolerance only when it threatens the conditions of toleration. The resolution leaves open how a given society is to judge whether a given view threatens those conditions. We propose a test with three factors: the capacity of the view's holders to act on it, the proximity of their action to the conditions it would impair, and the resilience of those conditions in the society in question. We argue that the threshold is relative to the society, that errors of over-restriction are themselves threats to the conditions of toleration and must be weighed as such, and that judgments of threat must be made by bodies whose findings can be contested in public.
\end{abstract}

\section{The problem}

Theorem~3.2 of POL-16-a shows that a policy satisfies the principles of preservation and necessity if and only if it tolerates a view exactly when the view does not threaten the conditions of toleration, and Corollary~3.3 shows that such a policy tolerates every view that is not active intolerance. The theorem is a characterization: it says which policy is right, given the facts about which views threaten. It does not say how those facts are to be known. \citet[\S35]{Rawls1971} held that the intolerant should be restricted only when the tolerant sincerely and with reason believe that their own security and that of the institutions of liberty are in danger, and \citet[ch.~7, n.~4]{Popper1945} that intolerance should be suppressed when it refuses rational argument and resorts to force. Both locate the threshold; neither supplies a way of applying it.

\section{A test}

Let $v$ be an active intolerant view, held in a society $S$. By Definition~2.3 of POL-16-a, the conditions of toleration are constitutional protection of conscience and expression, the rule of law, and a distribution of power that prevents any group from suppressing others.

\begin{definition}\label{d:threat}
The \emph{threat} that $v$ poses in $S$ is a function of three factors: (T1) \emph{capacity}, the ability of $v$'s holders to bring about the suppression $v$ prescribes, through numbers, organization, resources, or access to coercive means; (T2) \emph{proximity}, the degree to which their action, if unimpeded, would impair one of the conditions of toleration directly rather than through the free choices of others; (T3) \emph{resilience}, the capacity of $S$'s institutions to resist the impairment by ordinary legal means. Threat increases with T1 and T2 and decreases with T3.
\end{definition}

The factors explain familiar judgments. A small sect that preaches the suppression of unbelievers but has neither numbers nor arms scores low on T1. A movement that campaigns for an amendment abolishing religious liberty acts through the free choices of voters and scores low on T2; it is answered by argument and by the entrenchment of the liberty. An armed organization that intimidates judges or voters scores high on T1 and T2 at once. The American standard, under which advocacy may be restricted only when directed to inciting imminent lawless action and likely to produce it \citep{Brandenburg1969}, is a judgment of high T2 combined with a presumption of high T3.

\begin{proposition}\label{p:relative}
Whether a view threatens the conditions of toleration depends on the society in which it is held.
\end{proposition}

\begin{proof}
T1 and T3 are properties of the relation between the view's holders and the society's institutions. The same organization may have decisive capacity in a society whose institutions are new and fragile, and negligible capacity where they are old and entrenched. Hence the same view may meet the threshold in one society and fail it in another.
\end{proof}

The proposition explains why the militant democracy of post-war Germany, which permits the prohibition of parties that seek to abolish the constitutional order, and the more permissive American regime can both be defensible: they answer to different values of T3 \citep{Loewenstein1937,Kirshner2014}. It also means that a policy appropriate when institutions were fragile should be relaxed as they become resilient.

\begin{proposition}\label{p:overrestriction}
A restriction that is not required by the threshold itself impairs the conditions of toleration.
\end{proposition}

\begin{proof}
To restrict a view is to subject its expression or organization to legal sanction (Definition~2.2 of POL-16-a). A power to do so, exercised where the threshold is not met, is a power of some to suppress the views of others outside the constraint that justifies it; and the absence of such powers is one of the conditions of toleration. Over-restriction therefore damages the very conditions the policy exists to preserve.
\end{proof}

Proposition~\ref{p:overrestriction} shows that the costs of error are not one-sided. A judgment of threat must weigh the danger of tolerating a threatening view against the damage done by restricting one that does not threaten, and both are damage to the same conditions. When the evidence leaves it uncertain whether the threshold is met, the second danger is not to be discounted because it is the familiar one.

\section{Who judges}

The factors of Definition~\ref{d:threat} are matters of fact, and judgments of them can be biased by those who make them. A government that may restrict the views it judges threatening has an interest in finding its opponents threatening. Two requirements follow. Restriction should be imposed only by bodies independent of the parties to political competition, such as courts, under standards stated in advance. And their findings should be open to public contest. POL-20 concludes that a public sphere is possible only under institutional conditions that guarantee equal access to public debate; where those conditions hold, claims of threat can be challenged by those against whom they are made, and the judgment is less likely to track the interests of those who hold power.

\begin{objection}
Courts are slow, and a threat that has become imminent cannot wait for litigation.
\end{objection}

\begin{reply}
Ordinary criminal law already restricts violence, conspiracy and intimidation, and does so quickly. The test concerns the restriction of views and of organizations as such. A view whose holders pose an imminent threat will usually be acting in ways that the criminal law reaches already; restriction of the view itself is needed only where the threat lies in organization and capacity, which develop over time and can be judged with the deliberation the test requires.
\end{reply}

\begin{objection}
The factors are vague, and a vague test invites abuse.
\end{objection}

\begin{reply}
They are vague, and the requirements of independence and contestability are the response to that fact, not a denial of it. The alternative tests are no more precise: restricting all intolerant views is precise, but by Proposition~3.4 of POL-16-a it violates the principle of necessity whenever some intolerant view does not threaten.
\end{reply}

\section{Conclusion}

The threshold of restriction in POL-16 is met when an active intolerant view, held by people with the capacity to act on it, would directly impair the conditions of toleration in a society whose institutions could not resist by ordinary means. The threshold is relative to the society, the costs of error lie on both sides, and the judgment belongs to independent bodies whose findings are publicly contestable. So understood, the resolution of the paradox restricts intolerance only when it threatens the conditions of toleration, as the conclusion of POL-16 requires.

\begin{thebibliography}{9}\small
\bibitem[Brandenburg v. Ohio(1969)]{Brandenburg1969} \emph{Brandenburg v. Ohio}, 395 U.S. 444 (1969).
\bibitem[Kirshner(2014)]{Kirshner2014} Kirshner, A. S. (2014). \emph{A Theory of Militant Democracy: The Ethics of Combatting Political Extremism}. Yale University Press.
\bibitem[Loewenstein(1937)]{Loewenstein1937} Loewenstein, K. (1937). Militant democracy and fundamental rights, I. \emph{American Political Science Review}, 31(3), 417--432.
\bibitem[Popper(1945)]{Popper1945} Popper, K. R. (1945). \emph{The Open Society and Its Enemies}, vol.~1. Routledge.
\bibitem[Rawls(1971)]{Rawls1971} Rawls, J. (1971). \emph{A Theory of Justice}. Harvard University Press.
\end{thebibliography}

\end{document}
